\documentclass[12pt,a4paper]{article}

\usepackage{amsmath,amssymb,amsthm}
\usepackage{geometry}
\usepackage{setspace}
\usepackage{enumitem}
\usepackage[hidelinks]{hyperref}
\usepackage{bm}
\usepackage{booktabs}
\usepackage[T2A]{fontenc}
\usepackage[utf8]{inputenc}
\usepackage[russian]{babel}
\usepackage{caption}
\usepackage{graphicx}
\usepackage{float}
\usepackage{listings}
\usepackage{tikz}
\usetikzlibrary{shapes.geometric, arrows.meta, positioning}
\usepackage[table]{xcolor}
\lstset{
    basicstyle=\ttfamily\small,
    breaklines=true,
    escapeinside={(*@}{@*)}
}
\captionsetup{font=small}
\geometry{left=2cm,right=2.5cm,top=2.5cm,bottom=2.5cm}
\onehalfspacing

\newtheoremstyle{bigplain}
  {1em}
  {1em}
  {\normalfont\rightskip=0pt}
  {}
  {\bfseries\large}
  {.}
  {0.2em}
  {\thmname{#1}\thmnote{(\textbf{#3})}}

\theoremstyle{bigplain}
\newtheorem*{theorem}{Теорема}

\theoremstyle{plain}
\newtheorem{definition}{Определение}
\newtheorem{remark}{Замечание}










\title{Численные методы решения\\ трансцендентного уравнения одного вида}
\author{}
\date{}

\begin{document}
\maketitle

%====================================================================
\begin{abstract}
Работа посвящена поиску решений неавтономной системы
с релейным гистерезисом численными методами.
Рассматривается трансцендентное уравнение,
полученное в работе В.В. Евстафьевой при исследовании
 двухточечно-колебательных
решений системы.
Проанализирована Теорема~1 указанной работы,
в которой установлены  условия существования и единственности
решения с заданным периодом возврата.
Выполнен обзор и сравнительный анализ численных методов, применяемых для 
поиска корней нелинейных уравнений и 
реализованных в средах \textsc{Matlab} и \textsc{Julia}.
Для модельной трехмерной системы найдено численное решение трансцендентного уравнения,
определяющего момент первого переключения, и проведен сравнительный анализ с аналитическим решением.
\end{abstract}

\vspace{2\baselineskip}

\noindent\textbf{Предварительные сведения: }
Рассматривается система обыкновенных дифференциальных уравнений
\begin{equation}
\dot Y(t)=AY(t)+Bu(\sigma(t))+Kf(t),\qquad
\sigma(t)=(C,Y(t)),
\tag{1.1}
\label{eq:system}
\end{equation}
где $Y(t)\in\mathbb{R}^n$ — вектор состояния, постоянная матрица $A$ имеет простые ненулевые вещественные собственные числа,
векторы $B,C,K\in\mathbb{R}^n$ постоянны,
$f(t)$ — непрерывная ограниченная функция возмущения .
Нелинейность $u(\sigma)$ описывает двухпозиционное реле
с положительным гистерезисом, порогами $\ell_1<\ell_2$
и уровнями выхода $m_1<m_2$.\\

%====================================================================
\noindent\textbf{Приведем основные определения из работы } \cite{Eustafieva}
\vspace{-0.5\baselineskip}
%====================================================================



\begin{definition}
Точкой переключения называется состояние системы~(1.1), 
при котором входная функция $\sigma(t)$ достигает одного из пороговых значений 
$\ell_\mu$, $\mu=1,2$, а выходная функция $u(t)$ при этом меняет значение 
выхода $m_1$ на $m_2$ или наоборот.
\end{definition}

\begin{definition}
Гиперплоскостью переключения называется гиперплоскость
\[
L_\mu=\{Y\in\mathbb{R}^n:(C,Y)=\ell_\mu\},\qquad \mu=1,2.
\]
\end{definition}

\begin{definition}
Решение $Y(\cdot)$ системы~(1.1) называется
$T_r$-двухточечно-колебательным простейшего поведения,
если существуют вещественные положительные числа $\tau_1$ и $\tau_2$ такие, что
$\tau_1+\tau_2=T_r$, и точки $Y^{1}\in L_1$ и $Y^{2}\in L_2$,
для которых выполняются следующие условия:

\begin{enumerate}[label=\textbf{\arabic*)}, leftmargin=*]
\item
$Y(t_0+\kappa T_r)=Y^{1}$
\quad \text{и} \quad
$Y(t_0+\tau_1+\kappa T_r)=Y^{2},$
где $t_0\ge 0$ -- начальный момент времени,
$\kappa\in\mathbb{N}\cup\{0\}$;

\item
для всех $\kappa\in\mathbb{N}\cup\{0\}$ на полуинтервалах
$\Delta_\kappa^{1}=[t_0+\kappa T_r,\; t_0+\tau_1+\kappa T_r)$
имеет место равенство $u=m_1$, а на полуинтервале
$\Delta_\kappa^{2}=[t_0+\tau_1+\kappa T_r,\; t_0+(\kappa+1)T_r)$
имеет место равенство $u=m_2$.
\end{enumerate}
\end{definition}   




\noindent{Далее приведем вспомогательные сведения.}

%====================================================================
\subsection*{Формула Коши для неавтономной системы}
%====================================================================

Рассмотрим линейную систему
\begin{equation}
\dot Y(t)=AY(t)+Kf(t),
\tag{1.2}
\label{eq:linear}
\end{equation}
где $A$ — постоянная матрица, $f(t)$ — непрерывная функция.\\
Пусть
\[
Y(t)=e^{At}Z(t).
\]
Тогда
\[
\dot Y(t)=\frac{d}{dt}\bigl(e^{At}Z(t)\bigr)
= Ae^{At}Z(t)+e^{At}\dot Z(t).
\]
Из \eqref{eq:linear} получаем
\[
e^{At}\dot Z(t) = Kf(t) \implies \dot Z(t) = e^{-At}Kf(t).
\]
Интегрируя по времени от $t_0$ до $t$, имеем
\[
Z(t)-Z(t_0)=\int_{t_0}^{t} e^{-As}Kf(s)\,ds.
\]
Поскольку $Z(t_0)=e^{-At_0}Y(t_0)$, возвращаясь к $Y(t)$, получаем
\[
Y(t)
=
e^{A(t-t_0)}Y(t_0)
+
\int_{t_0}^{t} e^{A(t-s)}Kf(s)\,ds.
\]

\noindentВ общем случае для системы
\[
\dot Y(t)=AY(t)+g(t),
\]
где $g(t)$ — разрывно-непрерывная вектор-функция, формула Коши принимает вид
\begin{equation}
Y(t)
=
e^{A(t-t_0)}Y(t_0)
+
\int_{t_0}^{t} e^{A(t-s)}g(s)\,ds.
\tag{1.3}
\label{eq:Cauchy}
\end{equation}

\noindentВ системе \eqref{eq:system} роль внешнего воздействия играет функция
\[
g(t)=Bu(\sigma(t))+Kf(t).
\]

Формула \eqref{eq:Cauchy} является фундаментальным инструментом
для описания движения системы между моментами переключения,
когда значение $u(\sigma)$ постоянно и система становится линейной.


%====================================================================
\subsection*{Метод припасовывания. Вывод условий возврата}
%====================================================================

Рассматривается система
\begin{equation*}
\dot Y(t)=AY(t)+Bu(\sigma(t))+Kf(t),
\qquad
\sigma(t)=(C,Y(t)).
\end{equation*}

Пусть $Y^{1}\in L_1$, $Y^{2}\in L_2$ — точки переключения, $\tau_1$ - время перехода от $L_1$ до $L_2$ \\
$t_0+\tau_1$ — момент первого переключения,
$T_r=\tau_1+\tau_2$ — период возврата.

%--------------------------------------------------------------------
\paragraph{1. Движение на первом участке ($u=m_1$).}
%--------------------------------------------------------------------

При $t\in[t_0,t_0+\tau_1)$ имеем $u(t)=m_1$, поэтому система принимает вид
\[
\dot Y(t)=AY(t)+Bm_1+Kf(t).
\]

\noindentПо формуле Коши получаем
\begin{equation*}
Y(t)
=
e^{A(t-t_0)}Y^{1}
+
\int_{t_0}^{t}
e^{A(t-s)}(Bm_1+Kf(s))\,ds.
\end{equation*}

\noindentВ момент $t=t_0+\tau_1$ траектория достигает гиперплоскости $L_2$, поэтому
\begin{equation}
Y^{2}
=
e^{A\tau_1}Y^{1}
+
\int_{t_0}^{t_0+\tau_1}
e^{A(t_0+\tau_1-s)}
\bigl(Bm_1+Kf(s)\bigr)\,ds.
\tag{1.4}
\label{eq:Y2_expression}
\end{equation}

%--------------------------------------------------------------------
\paragraph{2. Движение на втором участке ($u=m_2$).}
%--------------------------------------------------------------------

При $t\in[t_0+\tau_1,t_0+T_r)$ имеем $u(t)=m_2$, и система принимает вид
\[
\dot Y(t)=AY(t)+Bm_2+Kf(t).
\]

\noindentПо формуле Коши:
\begin{equation}
Y(t)
=
e^{A(t-t_0-\tau_1)}Y^{2}
+
\int_{t_0+\tau_1}^{t}
e^{A(t-s)}(Bm_2+Kf(s))\,ds.
\tag{1.5}
\label{eq:second_arc}
\end{equation}

\noindentВ момент $t=t_0+T_r$ выполняется условие возврата:
\[
Y(t_0+T_r)=Y^{1}.
\]

\noindentПодставляя $t=t_0+T_r$ в \eqref{eq:second_arc}, получаем
\begin{equation}
Y^{1}
=
e^{A(T_r-\tau_1)}Y^{2}
+
\int_{t_0+\tau_1}^{t_0+T_r}
e^{A(t_0+T_r-s)}(Bm_2+Kf(s))\,ds.
\tag{1.6}
\label{eq:return_condition}
\end{equation}


%--------------------------------------------------------------------
\paragraph{3. Исключение $Y^{2}$ (припасовывание).}
%--------------------------------------------------------------------

Подставляя выражение \eqref{eq:Y2_expression} в
\eqref{eq:return_condition}, имеем
\begin{equation*}
Y^{1}=
e^{A(T_r-\tau_1)}
\left(
e^{A\tau_1}Y^{1}
+
\int_{t_0}^{t_0+\tau_1}
e^{A(t_0+\tau_1-s)}
\bigl(Bm_1+Kf(s)\bigr)\,ds
\right)
+
\int_{t_0+\tau_1}^{t_0+T_r}
e^{A(t_0+T_r-s)}
\bigl(Bm_2+Kf(s)\bigr)\,ds.
\end{equation*}

\noindent Используя свойство матричной экспоненты
$e^{A(T_r-\tau_1)}e^{A\tau_1}=e^{AT_r}$, получаем
\begin{equation*}
Y^{1}=
e^{AT_r}Y^{1}
+
\int_{t_0}^{t_0+\tau_1}
e^{A(t_0+T_r-s)}
\bigl(Bm_1+Kf(s)\bigr)\,ds
+
\int_{t_0+\tau_1}^{t_0+T_r}
e^{A(t_0+T_r-s)}
\bigl(Bm_2+Kf(s)\bigr)\,ds.
\end{equation*}

\noindent Перенося выражение $e^{AT_r}Y^{1}$ в левую часть, приходим к равенству
\begin{equation}
(E-e^{AT_r})Y^{1}
=
\int_{t_0}^{t_0+\tau_1}
e^{A(t_0+T_r-s)}
\bigl(Bm_1+Kf(s)\bigr)\,ds
+
\int_{t_0+\tau_1}^{t_0+T_r}
e^{A(t_0+T_r-s)}
\bigl(Bm_2+Kf(s)\bigr)\,ds.
\tag{1.7}
\label{eq:Y1_matching}
\end{equation}

\noindent При обратимости матрицы $E-e^{AT_r}$ из
\eqref{eq:Y1_matching} получаем выражение для точки переключения
\begin{equation*}
Y^{1}
=
(E-e^{AT_r})^{-1}
\left(
\int_{t_0}^{t_0+\tau_1}
e^{A(t_0+T_r-s)}
\bigl(Bm_1+Kf(s)\bigr)\,ds
+
\int_{t_0+\tau_1}^{t_0+T_r}
e^{A(t_0+T_r-s)}
\bigl(Bm_2+Kf(s)\bigr)\,ds
\right).
\end{equation*}

\noindent После замены
\begin{equation*}
Y^{1}
=
(E-e^{AT_r})^{-1}
\left(
\int_{0}^{\tau_1}
e^{A(T_r-\xi)}
\bigl(Bm_1+Kf(t_0+\xi)\bigr)\,d\xi
+
\int_{0}^{T_r-\tau_1}
e^{A\xi}
\bigl(Bm_2+Kf(t_0+T_r-\xi)\bigr)\,d\xi
\right).
\end{equation*}

\newpage
\noindent{\large\textbf{\mbox{Критерий существования решения с заданным периодом возврата}}}
\\
\noindent{\large\textbf{\mbox{и нулевым начальным моментом времени}}}

\begin{theorem} \cite{Eustafieva}
Существует единственное решение системы (1.1) с заданным периодом возврата тогда и только тогда, когда при некоторой фиксированной величине $T_r$ выполняются следующие условия:

\begin{enumerate}[label=\arabic*)]
\item уравнение
\begin{equation}
\left(
C,\,
(E-e^{AT_r})^{-1}e^{Ax}
\left(\int_0^x e^{-A\xi}(Bm_1+Kf(\xi))\,d\xi
+\int_0^{T_r-x} e^{A\xi}(Bm_2+Kf(T_r-\xi))\,d\xi
\right)
\right)=\ell_2
\tag{2.1}
\label{eq:condition21}
\end{equation}
имеет хотя бы одно решение, где $E$ -- единичная матрица, $x$ -- независимая переменная из интервала $(0,T_r)$;

\item для величины $\tau_1$, удовлетворяющей уравнению (2.1), выполняется равенство \newline
$(C,Y^1)=\ell_1$, где
\[
Y^1=(E-e^{AT_r})^{-1}\left(\int_0^{\tau_1} e^{A(T_r-\xi)}(Bm_1+Kf(\xi))\,d\xi+\int_0^{T_r-\tau_1} e^{A\xi}(Bm_2+Kf(T_r-\xi))\,d\xi\right),
\]
и интегралы
\[
K\int_0^{\tau_1} e^{-A\xi}f(\xi+kT_r)\,d\xi,
\qquad
K\int_0^{T_r-\tau_1} e^{A\xi}f((k+1)T_r-\xi)\,d\xi
\]
не зависят от $k\in\mathbb{N}\cup\{0\}$;

\item величина $\tau_1$, удовлетворяющая условию 2), является наименьшим (или единственным) решением уравнения
\[
\left(C,\,e^{Ax}Y^1+\int_0^x e^{A(x-\xi)}(Bm_1+Kf(\xi))\,d\xi
\right)=\ell_2,
\]
где $Y^1$ определена в условии 2), $x$ -- независимая переменная из интервала $(0,T_r)$;

\item величина $T_r-\tau_1$ является наименьшим (или единственным) решением уравнения
\[
\left(C,\,e^{Az}Y^2+\int_0^z e^{A\xi}(Bm_2+Kf(\tau_1+z-\xi))\,d\xi
\right)=\ell_1,
\]
где $\tau_1$ удовлетворяет условию 3), $z$ -- независимая переменная из интервала $(0,T_r)$,
\[
Y^2=e^{A\tau_1}Y^1+\int_0^{\tau_1} e^{A(\tau_1-\xi)}(Bm_1+Kf(\xi))\,d\xi,
\]
$Y^1$ определена в условии 2).
\end{enumerate}
\end{theorem}


\subsection*{Постановка задачи}

Согласно Теореме~1, существование $T_r$-двухточечно-колебательного решения
гарантируется, в первую очередь, существованием решения $\tau_1$
трансцендентного уравнения~\eqref{eq:condition21} (пункт~1),
а также выполнением дополнительных условий (пункты~2--4).

Напомним, что $\tau_1$ находится из условия,
что траектория системы,
построенная по формуле Коши и стартующая в момент времени $t_0=0$
из точки переключения $Y^{1}$,
достигает гиперплоскости переключения $L_2$
ровно в момент времени $t=\tau_1$.
Исключая координаты точек переключения
путём припасовывания решений на соседних участках движения,
получено уравнение \eqref{eq:condition21}
относительно $x\in(0,T_r)$. Пусть
\begin{equation*}
    Q_1 = \int_0^{x} e^{-A\xi}\bigl(Bm_1+Kf(\xi)\bigr)\,d\xi \qquad
    Q_2 = \int_0^{T_r-x} e^{A\xi}\bigl(Bm_2+Kf(T_r-\xi)\bigr)\,d\xi
\end{equation*}

\vspace{-15pt}
\noindent
\begin{minipage}[t]{0.65\textwidth}
\begin{equation}
\tag{3}\label{eq:trans_func}
F(x)=
\left(
C,\,
(E-e^{AT_r})^{-1} e^{Ax}
\left(
Q_1+Q_2
\right)
\right)
-\ell_2.
\end{equation}
\end{minipage}\hfill
\begin{minipage}[t]{0.34\textwidth}
\begin{equation}
\tag{3.1}\label{eq:trans_eq}
F(x)=0.
\end{equation}
\end{minipage}

\vspace{15pt}
Задача данной работы состоит в решении нелинейного (трансцендентного) уравнения $F(x)=0$ (\eqref{eq:trans_eq}), где $F(x)$ определена формулой \eqref{eq:trans_func}.

\vspace{5pt}
\noindent\textbf{Обзор численных методов решения нелинейных уравнений}

Методы решения нелинейных уравнений
обычно подразделяются на два основных класса: методы с локализацией корня и открытые методы~\cite{Press, Badr}.
В отличие от быстрых, но потенциально расходящихся методов второго
класса, методы первого гарантируют сходимость ценой снижения скорости.

Ниже схематично проиллюстрирована работа рассматриваемых в данной
работе методов поиска корня для скалярного уравнения $F(x)=0$.
(Для оценки скорости сходимости используется формула $|e_{k+1}| = C|e_k|^p$, где $e_k$ -- погрешность на $k$-й итерации, $C$ -- константа, $p$ -- порядок сходимости.)


\begin{itemize}

\item \textbf{Метод дихотомии (Bisection method, Bi).}

\noindent
\begin{minipage}{0.6\textwidth}
Метод дихотомии применяется для решения уравнения $F(x)=0$
при условии непрерывности функции $y=F(x)$ на $(a,b)$ и выполнении условия
F(a)F(b)<0,
что гарантирует существование корня.

На каждой итерации интервал делится пополам,
а новое приближение определяется формулой
\[
x_{k+1}=\frac{a_k+b_k}{2}, \quad k=0,1,2,\ldots, \quad a_0=A, b_0=B
\]
Далее выбирается подинтервал, содержащий корень.

Метод гарантированно сходится, однако обладает линейной скоростью сходимости (p = 1) \cite{Kalitkin,Press}.
\end{minipage}
\hfill
\begin{minipage}{0.3\textwidth}
\centering
\vspace{-45pt}
\includegraphics[width=\textwidth]{Bisection.png}

\small Рис. 1 Bisection
\end{minipage}


\item \textbf{Метод деления отрезка на три части (Trisection method, Tri).}

\noindent
\begin{minipage}{0.6\textwidth}
Метод трисекции является модификацией метода
дихотомии. В данном случае интервал
$(a,b)$ делится на три равные части.
Вычисляются значения функции в точках
\[
x_{k+1}^1=a+\frac{b_k-a_k}{3},
\quad
x_{k+1}^2=a+\frac{2(b_k-a_k)}{3},
\]
\vspace{-5pt}
\[
k=0,1,2,\ldots, \quad a_0=A, b_0=B
\]
после чего выбирается подинтервал,
в котором происходит смена знака функции.
Метод гарантированно сходится,
но обладает линейной скоростью сходимости (p = 1) и требует большего числа вычислений
значений функции на каждой итерации.
\end{minipage}
\hfill
\begin{minipage}{0.3\textwidth}
\centering
\vspace{0pt}
\includegraphics[width=\textwidth]{Trisection.png}

\small Рис. 2 Trisection
\end{minipage}

\item \textbf{Метод ложного положения
(False Position method, FP).}

\noindent
\begin{minipage}{0.6\textwidth}
Метод ложного положения использует
линейную интерполяцию функции
на интервале $(a,b)$.
Новое приближение определяется
точкой пересечения секущей,
проходящей через точки
$(a,F(a))$ и $(b,F(b))$,
с осью абсцисс:
\[
x=b-\frac{F(b)(b-a)}{F(b)-F(a)}, \quad k=0,1,2,\ldots, \quad a_0=A, b_0=B
\]
Метод гарантированно сходится, но обладает линейной скоростью сходимости (p = 1).
Также данный метод сохраняет интервал локализации корня. \cite{Press}
\end{minipage}
\hspace{7pt}
\hfill
\begin{minipage}{0.38\textwidth}
\centering
\vspace{0pt}
\includegraphics[width=\textwidth]{FPM.png}

\small Рис. 3 FP
\end{minipage}


\item \textbf{Метод Ньютона
(Newton--Raphson method, NR).}

\noindent
\begin{minipage}{0.6\textwidth}
Метод Ньютона относится к открытым методам
и основан на аппроксимации функции
касательной в текущей точке.
Итерационная формула имеет вид
\[
x_{k+1}=x_k-\frac{F(x_k)}{F'(x_k)}.
\]
Метод обладает квадратичной
скоростью сходимости (p = 2),
однако требует вычисления
производной функции
и чувствителен к выбору
начального приближения (точка A на рис. 4) \cite{Kalitkin,Press}.
\end{minipage}
\hfill
\begin{minipage}{0.2\textwidth}
\centering
\vspace{-10pt}
\includegraphics[width=\textwidth]{Newton.png}

\small Рис. 4 NR
\end{minipage}
\\

\item \textbf{Метод секущих (Secant method, Se).}

\noindent
\begin{minipage}{0.6\textwidth}
Метод секущих представляет собой
модификацию метода Ньютона,
в которой производная функции
заменяется разделенной разностью:
\[
x_{k+1}=x_k-
\frac{x_k-x_{k-1}}
{F(x_k)-F(x_{k-1})}F(x_k).
\]
Метод не требует вычисления производной
функции и обладает сверхлинейной
скоростью сходимости \\ (p $\approx 1.618$),
однако не гарантирует
сходимости итерационного процесса \cite{Press}.
\end{minipage}
\hfill
\begin{minipage}{0.2\textwidth}
\centering
\vspace{-10pt}
\includegraphics[width=\textwidth]{Sectant.png}

\small Рис. 5 Sectant
\end{minipage}

\item \textbf{Модифицированный метод секущих
(Modified Secant method, MSe).}

В данном методе производная
заменяется разностью значений
функции в точках $x$ и $x+\delta x$,
что позволяет использовать
одно начальное приближение.
Итерационная формула имеет вид
\[
x_{k+1}=x_k-
\frac{\delta x_k F(x_k)}
{F(x_k+\delta x_k)-F(x_k)}, \quad \text{где} \quad \delta x_k \text{- малое приращение аргумента}
\]

Порядок сходимости данного метода зависит
от стратегии выбора $\delta x_k$: в случае фиксации
шага ($\delta x_k = \text{const}$) наблюдается линейная сходимость, 
тогда как выбор приращения, пропорционального значению 
самой функции ($\delta x_k = F(x_k)$), обеспечивает
квадратичную сходимость, делая алгоритм эквивалентным
методу Стеффенсена \cite{Press,Steffensen}.

\item \textbf{Гибридный метод ложной позиции и модифицированного метода секущих
(False Position--Modified Secant method, FP--MSe).}

Метод FP--MSe объединяет преимущества метода ложной позиции
и открытого модифицированного метода секущих.
Пусть на $k$-й итерации задан интервал $(a_k,b_k)$,
на концах которого $F(a_k)F(b_k)<0$.

Сначала вычисляется приближение методом ложной позиции:
\[
x_k^{(FP)} = a_k - \frac{F(a_k)(b_k-a_k)}{F(b_k)-F(a_k)}.
\]
Затем это приближение уточняется по схеме
модифицированного метода секущих:
\[
x_k^{(MSe)} = x_k^{(FP)} -
\frac{\delta x_k\,F(x_k^{(FP)})}
{F(x_k^{(FP)}+\delta x_k)-F(x_k^{(FP)})}.
\]

Итоговое новое приближение выбирается по правилу
\[
x_{k+1}=
\begin{cases}
x_k^{(MSe)}, & \text{если } x_k^{(MSe)}\in(a_k,b_k)
\text{ и } |F(x_k^{(MSe)})|<|F(x_k^{(FP)})|,\\[2mm]
x_k^{(FP)}, & \text{в противном случае}.
\end{cases}
\]
После этого интервал локализации корня обновляется
по знаку функции $F(x_{k+1})$.
Данный метод обладает гарантированной линейной сходимостью \cite{Badr}..

\item \textbf{Метод Галлея (Halley method).}

Метод Галлея является обобщением метода Ньютона,
поскольку использует информацию о второй производной функции.
Он получается при учёте второго члена разложения
функции $F(x)$ в ряд Тейлора в окрестности текущего
приближения~$x_k$.

Итерационная формула имеет вид
\[
x_{k+1}=x_k-
\frac{F(x_k)}{F'(x_k)}
\left(
1-\frac{F(x_k)F''(x_k)}{2F'(x_k)^2}
\right)^{-1}.
\]

Метод Галлея обладает кубической скоростью
сходимости.
Однако применение метода оправдано только
в тех случаях, когда вычисление второй
производной $F''(x)$ не приводит к
существенным дополнительным затратам
\cite{Press}.


\item \textbf{Метод Риддера (Ridders method, Rid).}

Метод Риддера является модификацией метода ложной позиции
и основан на экспоненциальной аппроксимации функции
на интервале $(x_1,x_2)$, где $F(x_1)F(x_2)<0$.

Идея заключается в том, что на интервале $(x_1, x_2)$ 
вводится вспомогательная функция $h(x) = F(x)e^{ax}$, 
где параметр $a$ подбирается таким образом, чтобы экспоненциальный 
множитель компенсировал кривизну исходной функции $F(x)$. 
В результате этого узловые точки $(x_1, h(x_1))$, $(x_3, h(x_3))$ 
и $(x_2, h(x_2))$ оказываются на одной прямой, 
что позволяет эффективно применить метод ложной позиции.

Сначала вычисляется значение функции
в средней точке
\[
x_3=\frac{x_1+x_2}{2}.
\]

Далее вводится вспомогательная функция $h(x)=F(x)e^{ax}$,

где параметр $a$ выбирается таким образом,
чтобы выполнялось условие
\[
h(x_3)=\frac{h(x_1)+h(x_2)}{2} \implies 2F(x_3)e^{ax_3}=F(x_1)e^{ax_1}+F(x_2)e^{ax_2}
\]
\[
\implies 2F(x_3)e^{a(x_3-x_1)}=F(x_1)+F(x_2)e^{a(x_2-x_1)} \implies F(x_1)-2F(x_3)e^{Q}+F(x_2)e^{2Q}=0,
\]

где $Q=a(x_3-x_1)$. \\
Решая это уравнение относительно $e^{Q}$, получаем:
\begin{equation}
e^{Q}
=
\frac{F(x_3)+\operatorname{sign}(F(x_2))
\sqrt{F(x_3)^2-F(x_1)F(x_2)}}
{F(x_2)}.
\tag{4.1}
\label{eq:Ridders_e_Q}
\end{equation}

Для нахождения нового приближения к корню $x_4$ применяется метод
ложной позиции к преобразованной функции $h(x) = F(x)e^{ax}$. 
Итерационный шаг записывается в виде:
\[
x_4 = x_3 - h(x_3) \frac{x_3 - x_1}{h(x_3) - h(x_1)},
\]
где $h(x_1) = F(x_1)e^{ax_1}$, $h(x_3) = F(x_3)e^{ax_3} = F(x_3)e^{ax_1}e^Q$.
\[
x_4 = x_3 - (x_3 - x_1) \frac{F(x_3)e^Q}{F(x_3)e^Q - F(x_1)}.
\]

Подставляя в данное выражение вместо $e^Q$ правую 
часть \eqref{eq:Ridders_e_Q}, получаем формулу нового приближения:
\[
x_4 = x_3 + (x_3 - x_1)
\frac{\operatorname{sign}(F(x_1) - F(x_2))\,F(x_3)}
{\sqrt{F(x_3)^2 - F(x_1)F(x_2)}}.
\]

После вычисления $x_4$ выбирается один из подинтервалов $(x_1,x_4)$ или $(x_4,x_2)$,
на концах которого функция применяет значения разных знаков.

Метод обладает квадратичной скоростью сходимости~\cite{Press}.




\item \textbf{Метод Брента
(Van Wijngaarden--Dekker--Brent method).}

Метод Брента является гибридным методом,
объединяющим методы дихотомии, секущих
и обратной квадратичной интерполяции.

В процессе работы алгоритм непрерывно обновляет три ключевые точки:
\begin{itemize}
    \item $b$ — наилучшее текущее приближение к корню (модуль функции минимален);
    \item $c$ — точка, гарантирующая локализацию корня (выполняется $F(b)F(c) < 0$);
    \item $a$ — значение наилучшего приближения $b$ с предыдущей итерации.
\end{itemize}

Когда значения функции в точках $a, b, c$ попарно различны,
метод Брента применяет обратную квадратичную интерполяцию. 

Запишем интерполяционный многочлен Лагранжа для трех узлов $(F(a), a), (F(b), b), (F(c), c)$:
\begin{equation*}
\begin{aligned}
    x(y) &= \frac{(y - F(b))(y - F(c))}{(F(a) - F(b))(F(a) - F(c))} a 
         &+ \frac{(y - F(a))(y - F(c))}{(F(b) - F(a))(F(b) - F(c))} b \\
         &+ \frac{(y - F(a))(y - F(b))}{(F(c) - F(a))(F(c) - F(b))} c
\end{aligned}
\end{equation*}

Подставляем $y = 0$:
\begin{equation*}
    x = \frac{F(b)F(c) a}{(F(a) - F(b))(F(a) - F(c))} + \frac{F(a)F(c) b}{(F(b) - F(a))(F(b) - F(c))} + \frac{F(a)F(b) c}{(F(c) - F(a))(F(c) - F(b))}
\end{equation*}

Пусть $R = \frac{F(b)}{F(c)}$, $S = \frac{F(b)}{F(a)}$ и $T = \frac{F(a)}{F(c)}$:

Тогда предыдущая формула может быть переписана в виде
\begin{equation*}
    x = b + \frac{P}{Q}, \quad \text{где}
\end{equation*}
\begin{equation*}
    P = S \cdot \big[ T(R - T)(c - b) - (1 - R)(b - a) \big]
    \qquad Q = (T - 1)(R - 1)(S - 1)
\end{equation*}

Метод Брента применяет следующую стратегию выбора поправки $d = \frac{P}{Q}$:
\begin{itemize}
    \item \textbf{Обратная квадратичная интерполяция}: используется, если все три текущих значения $F(a), F(b), F(c)$ попарно различны.
    \item \textbf{Метод секущих}: применяется, если $F(a) = F(c)$ (ситуация $F(b)=F(c)$ невозможна из-за условия локализации $F(b)F(c)<0$, а при $F(a)=F(b)$ возникает деление на ноль, что немедленно запускает шаг дихотомии). Формула упрощается до линейной интерполяции: $P = S(c - b)$, $Q = S - 1$.
    \item \textbf{Метод дихотомии}: применяется, если новое приближение $b + d$ выходит за его границы, интерполяционный шаг убывает менее чем в два раза по сравнению с прошлыми итерациями, или достигнут предел заданной точности.
\end{itemize}

Метод сходится сверхлинейно ($p \approx 1.84$) на гладких
функциях, гарантируя в худшем случае линейную скорость сходимости
\cite{Press}.



\item \textbf{Методы типа Стеффенсена (Steffensen-like methods).}

Классический метод Стеффенсена~\cite{Steffensen} является итерационным
методом решения нелинейных уравнений $F(x)=0$, не требующим вычисления
производной функции. Итерационная схема в классическом случае имеет вид:
\begin{equation}
    x_{k+1} = x_k - \frac{F(x_k)^2}{F(x_k + F(x_k)) - F(x_k)}, \quad k = 0, 1, 2, \dots
    \tag{4.2}
    \label{eq:original_steffensen}
\end{equation}
Этот метод обладает квадратичной скоростью сходимости.

В работе~\cite{Traub} Джозеф Трауб обобщил данный метод, введя варьируемый параметр~$\gamma_k \neq 0$:
\[
x_{k+1} = x_k - \frac{\gamma_k F(x_k)^2}{F(x_k + \gamma_k F(x_k)) - F(x_k)}.
\]
Аппроксимируя производную $F'(x_k)$ через разделенную
разность $F[x_k, x_{k-1}]$, Трауб получил следующую формулу для
параметра $\gamma_k$:
\[
\gamma_k = -\frac{1}{N_1'(x_k)} = - \frac{1}{F[x_k, x_{k-1}]} = - \frac{x_k - x_{k-1}}{F(x_k) - F(x_{k-1})}, \quad k \ge 1,
\]
где $N_1(t)$ — интерполяционный многочлен Ньютона первой степени (общая формула: $N_n(t) = F(x_0) + \sum_{k=1}^{n} F[x_0, x_1, \dots, x_k] \prod_{j=0}^{k-1} (t - x_j)$).

Этот подход повысил порядок сходимости до сверхквадратичного: $1 + \sqrt{2} \approx 2.414$.

В работе \cite{Dzunic} идея использования памяти была развита за счет 
применения более точной аппроксимации производной $F'(x_k)$. Авторы
предложили интерполировать функцию $F(x)$ многочленом Ньютона второй
степени $N_2(t)$. Этот многочлен строится по трем доступным узлам
из текущей и предыдущей итераций: $x_k$, $x_{k-1}$ и $w_{k-1}$. Узел
$w_{k-1}$ вычисляется на основе данных предыдущего шага как $w_{k-1} = x_{k-1} + \gamma_{k-1} F(x_{k-1})$.

Многочлен Ньютона, интерполирующий функцию $F(x)$ в этих узлах, записывается
через разделенные разности следующим образом:
\[
N_2(t) = F(x_k) + F[x_k, x_{k-1}](t - x_k) + F[x_k, x_{k-1}, w_{k-1}](t - x_k)(t - x_{k-1}).
\]
Для получения аппроксимации производной $F'(x_k)$ этот многочлен
дифференцируется по переменной $t$, и результат вычисляется в точке текущего приближения~$t = x_k$:
\[
N_2'(x_k) = \left. \frac{d}{dt} N_2(t) \right|_{t=x_k} = F[x_k, x_{k-1}] + F[x_k, x_{k-1}, w_{k-1}](x_k - x_{k-1}).
\]
Таким образом, модифицированная итерационная схема с параметром $\gamma_k$ принимает вид:
\begin{equation}
x_{k+1} = x_k - \frac{\gamma_k\,F(x_k)^2}{F\!\bigl(x_k+\gamma_k F(x_k)\bigr)-F(x_k)},
\qquad
\begin{aligned}
\gamma_0 &\text{ задаётся начальным условием},\\
\gamma_k &= -\frac{1}{N_2'(x_k)}, \quad k\ge 1.
\end{aligned}
\tag{4.3}
\label{eq:modified_steffensen}
\end{equation}
Предложенный метод обладает кубической скоростью сходимости.

В данной работе из этого раздела будут использованы формулы \eqref{eq:original_steffensen} (далее St или Steffensen) и \eqref{eq:modified_steffensen} (далее MSt или Modified Steffensen).
\end{itemize}





\noindent\textbf{Реализация численных методов в \textsc{Matlab} и \textsc{Julia}.}

Для численного решения скалярного уравнения $F(x)=0$
в среде \textsc{Matlab} может быть использована функция \texttt{fzero} \cite{Matlab}.
Она предназначена для поиска корня скалярной функции
и реализует гибридную схему, основанную на комбинации
методов дихотомии, секущих и обратной квадратичной интерполяции.

\medskip

\noindent
В среде \textsc{Matlab} базовый вызов функции поиска корня имеет вид:
\[
x = \texttt{fzero(@F,\,init)}
\]

\noindent
\texttt{x}\hspace{1.8em}--- найденный корень уравнения $F(x)=0$;\\
\texttt{@F}\hspace{1em}--- указатель на скалярную функцию $F(x)$;\\
\texttt{init}\hspace{0.6em}--- начальное приближение или интервал.

\medskip

\noindent
Для получения дополнительной информации о ходе вычислений и управления параметрами доступен расширенный вызов:
\[
\texttt{[x, fval, exitflag, output] = fzero(@F,\,init,\,options)}
\]

\noindent
\texttt{fval} --- значение функции в найденной точке; \\
\texttt{exitflag} --- код причины остановки алгоритма; \\
\texttt{output} --- структура, содержащая подробную информацию о процессе решения; \\
\texttt{options} --- структура, определяющая дополнительные параметры поиска корня (например, задание критериев остановки, настройка вывода информации по итерациям);

\medskip

\noindent\textbf{Дополнительные инструменты.}

\noindent
Помимо \texttt{fzero}, специализированные задачи поиска корней в \textsc{Matlab} решаются следующими встроенными функциями:
\begin{itemize}
    \vspace{-2pt}
    \item \texttt{roots(p)}~--- вычисление всех (включая комплексные) корней алгебраического многочлена, заданного вектором коэффициентов \texttt{p} \cite{Matlab};
    \vspace{-2pt}
    \item \texttt{fsolve(fun, x0)}~--- решение многомерных систем нелинейных уравнений \cite{Matlab};
    \vspace{-2pt}
    \item \texttt{solve(eqn,\,var)} --- функция решает \texttt{eqn} относительно \texttt{var} \cite{Matlab}.\\
    \texttt{eqn} --- символьное уравнение или система уравнений;\\
    \texttt{var} --- переменная, относительно которой решается уравнение.
    \vspace{-2pt}
    \item \texttt{vpasolve(eqn,\,var,\,init)} --- находит численное решение \texttt{eqn} относительно \texttt{var} \cite{Matlab}.\\
    \texttt{init} --- начальное приближение или интервал поиска.
\end{itemize}

Метод \texttt{fzero} в данной работе используется при составлении таблицы 1 (см. далее). 
Дополнительные инструменты приведены для справки.
Более подробную информацию можно найти в документации \cite{Matlab}.

\medskip

В языке \textsc{Julia} численное решение уравнения $F(x)=0$
реализуется с помощью пакета \texttt{Roots.jl}, основной
интерфейс которого задаётся функцией \texttt{find\_zero}~\cite{Julia}:
\[
x = \texttt{find\_zero(F,\,init,\,method)}.
\]
\\

\vspace{-1cm}
\begin{samepage}
\noindent
\texttt{F}\hspace{1.8em}--- функция;\\
\texttt{init}\hspace{1em}--- начальное приближение или интервал;\\
\texttt{method}\hspace{0.2em}--- используемый численный метод (например, \texttt{Bisection()}, \texttt{FalsePosition()}, \texttt{Secant()}, \texttt{Roots.Newton()}, \texttt{Roots.Halley()}, \texttt{Roots.Brent()}, \texttt{Roots.Ridders()}, \texttt{Order2()}).
\end{samepage}
\medskip

Таким образом, \textsc{Matlab} и \textsc{Julia}
предоставляют средства для численного решения
трансцендентного уравнения~\eqref{eq:trans_eq},
позволяя использовать или реализовать различные
численные методы поиска корней.
\\




\noindent\textbf{Пример решения трансцендентного уравнения.}

В настоящем разделе уравнение~\eqref{eq:trans_eq}
решается численными методами
для модельной системы, приведенной в примере статьи
В.\,В.~Евстафьевой \cite{Eustafieva}.
Рассматривается система порядка $n=3$
со следующими матрицами и параметрами:
\[
A=
\begin{pmatrix}
-32 & -18.75 & -22.5\\
37 & 21 & 29\\
7 & 4.25 & 4.5
\end{pmatrix},
\quad
B=
\begin{pmatrix}
-2.5\\ 4\\ 0.5
\end{pmatrix},
\quad
C=
\begin{pmatrix}
-3\\ -1.5\\ -6
\end{pmatrix},
\quad
K=
\begin{pmatrix}
1.5\\ -2\\ -0.5
\end{pmatrix},
\]
\[
m_1=-5,\qquad
m_2=3.842556,\qquad
\ell_1=-6,\qquad
\ell_2=6.629453,\qquad
T_r= \frac{3\pi}{2},
\]
\[
f(t)=0.5+\sin(2t+1.815774)+3\sin(4t+1.62).
\]

После подстановки параметров в~\eqref{eq:trans_func} получаем
\vspace{-9pt}
\begin{equation}
\begin{aligned}
F(x) = \Biggl( 
    &\begin{pmatrix} -3 \\ -1.5 \\ -6 \end{pmatrix}^{\!\!T}, 
    \left(E-e^{A\frac{3\pi}{2}}\right)^{-1} e^{Ax} 
    \biggl(\int_0^{x} e^{-A\xi}\left(
            -5 \begin{pmatrix} -2.5 \\ 4 \\ 0.5 \end{pmatrix} 
            + \begin{pmatrix} 1.5 \\ -2 \\ -0.5 \end{pmatrix} f(\xi) 
        \right) d\xi 
        \\[-2mm]
        &\hspace{-2.5em}
        + \int_0^{\frac{3\pi}{2}-x} e^{A\xi} 
        \left( 
            3.842556 \begin{pmatrix} -2.5 \\ 4 \\ 0.5 \end{pmatrix} 
            + \begin{pmatrix} 1.5 \\ -2 \\ -0.5 \end{pmatrix} f\left(\frac{3\pi}{2}-\xi\right) 
        \right) d\xi 
    \biggr) 
\Biggr) - 6.629453
\end{aligned}
\tag{5}
\label{eq:trans_func_specific}
\end{equation}

Решение уравнения \eqref{eq:trans_func_specific} выполнено численно на интервале $(0, \frac{3\pi}{2})$ с точностью $10^{-12}$, которая обеспечивается критериями остановки, принятыми для каждого метода (невязка и/или длина интервала).
Вычисления выполнены в среде \textsc{Julia} и \textsc{Matlab}.
Ниже приведен график функции \eqref{eq:trans_func} на интервале $(0, \frac{3\pi}{2})$.

\newpage

\setcounter{figure}{5}
\begin{figure}[h]
\centering
\begin{minipage}{0.82\textwidth}
    \centering
    \includegraphics[width=\linewidth]{F_x.png}
    \caption{График функции $F(x)$ (условие переключения)}
\end{minipage}
\end{figure}
\vspace{-12pt}
Данный график построен при помощи библиотеки \texttt{matplotlib} в \texttt{Python} \cite{matplotlib}.

\vspace{4pt}
Ниже приведем код и блок-схему метода MSe с подробным разбором. Реализация остальных методов приведена в приложении (\texttt{numMethods.jl}).



\vspace{-10pt}
\begin{lstlisting}[
    basicstyle=\ttfamily\bfseries\small,
    aboveskip=0pt,
    belowskip=0pt
]

function modified_secant(F::Function, x0::Float64; tol, maxiter)
    # F - (*@\text{функция}@*), x0 - (*@\text{начальное приближение}@*), tol - (*@\text{требуемая точность}@*), maxiter - (*@\text{предельное число итераций}@*).
    x = x0 # (*@\text{Начальное приближение}@*)
    it = 0 # (*@\text{Счётчик итераций}@*)
    while it < maxiter 
        fx = F(x)  # (*@\text{Вычисление значения функции в текущей точке}@*)
        # (*@\text{Проверка критерия остановки по невязке}@*)
        abs(fx) < tol && return x, it 
        # (*@\text{Шаг пропорционален значению функции (эквивалент метода Стеффенсена)}@*)
        d = fx 
        denom = F(x + d) - fx # (*@\text{Разность значений функции}@*)
        @assert denom != 0.0 "(*@\text{Деление на ноль!}@*)"
        xnew = x - d * fx / denom # (*@\text{Вычисление нового приближения}@*)
        # (*@\text{Проверка критерия остановки по приращению аргумента}@*)
        if abs(xnew - x) < tol
            return xnew, it + 1
        end
        # (*@\text{Обновление текущего приближения и счётчика итераций}@*)
        x = xnew
        it += 1
    end
    return x, it
end

\end{lstlisting}


\begin{figure}[h!]
\centering
\begin{tikzpicture}[
    node distance=1.5cm and 2.5cm,
    startstop/.style={rectangle, rounded corners, minimum width=3cm, minimum height=1cm, text centered, draw=black, fill=gray!10},
    process/.style={rectangle, minimum width=4cm, minimum height=1cm, text centered, draw=black},
    decision/.style={diamond, aspect=2, minimum width=3cm, minimum height=1cm, text centered, draw=black},
    arrow/.style={thick, -{Stealth[scale=1.2]}}
    ]

    % Узлы
    \node (start) [startstop] {Начало};
    
    \node (init) [process, below=0.8cm of start] {$x = x_0$, $it = 0$};
    
    \node (loop) [decision, below=0.8cm of init] {$it < maxiter$?};
    
    \node (fx) [process, below=0.8cm of loop] {$fx = F(x)$};
    
    \node (tol1) [decision, below=0.8cm of fx] {$|fx| < tol$?};
    
    \node (calc) [process, below=0.8cm of tol1] {
        \begin{tabular}{c}
        $d = fx$ \\ 
        $denom = F(x + d) - fx$ 
        \end{tabular}
    };
    
    \node (assert) [decision, below=0.8cm of calc] {$denom \neq 0$?};
    
    \node (xnew) [process, below=0.8cm of assert] {$x_{new} = x - d \cdot \frac{fx}{denom}$};
    
    \node (tol2) [decision, below=0.8cm of xnew] {$|x_{new} - x| < tol$?};
    
    \node (update) [process, below=0.8cm of tol2] {$x = x_{new}$, $it = it + 1$};
    
    % Блоки возврата
    \node (ret1) [startstop, right=2.5cm of tol1] {Возврат $x, it$};
    \node (ret2) [startstop, right=2.5cm of assert] {Ошибка: Деление на 0};
    \node (ret3) [startstop, right=2.5cm of tol2] {Возврат $x_{new}, it+1$};
    \node (ret4) [startstop, right=3.5cm of loop] {Возврат $x, it$};

    % Связи (Стрелки)
    \draw [arrow] (start) -- (init);
    \draw [arrow] (init) -- (loop);
    
    \draw [arrow] (loop) -- node[anchor=east] {Да} (fx);
    \draw [arrow] (loop) -- node[anchor=south] {Нет} (ret4);
    
    \draw [arrow] (fx) -- (tol1);
    
    \draw [arrow] (tol1) -- node[anchor=east] {Нет} (calc);
    \draw [arrow] (tol1) -- node[anchor=south] {Да} (ret1);
    
    \draw [arrow] (calc) -- (assert);
    
    \draw [arrow] (assert) -- node[anchor=east] {Да} (xnew);
    \draw [arrow] (assert) -- node[anchor=south] {Нет} (ret2);
    
    \draw [arrow] (xnew) -- (tol2);
    
    \draw [arrow] (tol2) -- node[anchor=east] {Нет} (update);
    \draw [arrow] (tol2) -- node[anchor=south] {Да} (ret3);
    
    % Возврат цикла
    \draw [arrow] (update.west) -- ++(-1.5,0) |- (loop.west);

\end{tikzpicture}
\caption{Блок-схема алгоритма модифицированного метода секущих (MSe)}
\label{fig:mse_flowchart}
\end{figure}

В представленной блок-схеме используются математические формулы, соответствующие приведенному ранее описанию алгоритма модифицированного метода секущих.




\newpage



\noindent\textbf{Применение алгоритма к заданному примеру.}

Для решения рассматриваемого трансцендентного уравнения $F(x) = 0$ 
(где $F(x)$ задана формулой \eqref{eq:trans_func_specific})
метод вызывается со следующими параметрами:
начальное приближение задано как $x_0 = 1.5$, 
требуемая точность $tol = 10^{-12}$,
максимальное число итераций $maxiter = 100$. 

Ниже представлен фрагмент кода, демонстрирующий вызов функции
для конкретных значений:

\begin{lstlisting}[basicstyle=\ttfamily\bfseries\small]
# (*@\text{Вызов функции поиска корня}@*)
root, iterations = modified_secant(F, x0, tol, maxiter)
\end{lstlisting}

В результате выполнения программы для модельной системы алгоритм
успешно сходится. Искомое значение достигнуто на \textbf{N-й}
итерации и составляет \textbf{$x = \dots$}. \\

\vspace{-15pt}
Численные значения, представленные в таблице~\ref{tab:tau_compare}, получены на основе кода из приложения (\texttt{numMethods.jl}).

\vspace{10pt}
\begin{table}[htbp]
\centering
\small
\setlength{\tabcolsep}{6pt}
\renewcommand{\arraystretch}{1.15}



\rowcolors{1}{gray!15}{white}

\begin{tabular}{@{}lccc@{}}
\toprule
\rowcolor{gray!30}
Method & $\tau_1$ & Number of iterations & $|\tau_1 - \pi/2|$ \\
\midrule
Bisection         & 1.570799848595700 & 42 & $3.5218\cdot10^{-6}$ \\
Trisection        & 1.570799848553825 & 27 & $3.5218\cdot10^{-6}$ \\
False Position    & 1.551712098307423 & 200 & $1.9084\cdot10^{-6}$ \\
Newton--Raphson   & 1.570799848837137 & 57  & $3.5220\cdot10^{-6}$ \\
Secant            & 1.570799848580684 & 16  & $3.5218\cdot10^{-6}$ \\
Modified Secant   & 1.570799848589167 & 5 & $3.5218\cdot10^{-6}$ \\
FP--MSe           & 1.570799848900535 & 7  & $3.5221\cdot10^{-6}$ \\
Halley            & 1.570799848405245 & 5  & $3.5216\cdot10^{-6}$ \\
Ridders           & 1.570799848862333 & 11  & $3.5221\cdot10^{-6}$ \\
VW--Brent         & 1.570799848737147 & 19 & $3.5219\cdot10^{-6}$ \\
Steffensen        & 1.570799848924731 & 11  & $3.5221\cdot10^{-6}$ \\
\midrule
Analytical value  & $\pi/2 \approx 1.570796$ & -- & -- \\
\bottomrule
\end{tabular}

\caption{Сравнение численных методов решения уравнения~\eqref{eq:trans_eq}}
\label{tab:tau_compare}

\end{table}

Из таблицы видно, что все рассматриваемые методы
дают практически одно и то же значение момента
первого переключения:
\[
\tau_1 \approx 1.57079985.
\]
Отклонение от аналитической оценки
$\tau_1 \approx \pi/2$
составляет порядка $3.5\cdot 10^{-6}$.

Следует отметить, что приведённые в таблице
результаты соответствуют случаю достаточно близкого
начального приближения к искомому корню.
В этой ситуации методы MSe, FP--MSe, Halley, Ridders и Steffensen
демонстрируют быструю сходимость.


Для исследования влияния начальных данных на скорость сходимости
численных методов была построена карта сходимости на плоскости
$(x_0,x_1)$.
Параметры $x_0$ и $x_1$ задают начальные данные алгоритмов поиска корня.
Для интервальных методов (Bisection, False Position, Trisection,
Ridders, Brent и FP--MSe) величины
$x_0$ и $x_1$ определяют границы начального интервала.
Для метода Secant они используются как два начальных приближения.
Для методов Newton, Halley, MSe и метода
Steffensen в качестве начального приближения принималась середина
интервала
\[
x_m=\frac{x_0+x_1}{2}.
\]
Для каждой пары $(x_0,x_1)$ выполнялся запуск всех рассматриваемых
методов, после чего фиксировался метод, достигший решения за
наименьшее число итераций. Полученные результаты представлены на
(рис.~8).


\newpage

\begin{figure}[h]
\centering
\begin{minipage}{1.1\textwidth}
    \centering
    \includegraphics[width=\linewidth]{map.png}
    \small Рис. 7 Карта сходимости методов
\end{minipage}
\end{figure}

\newpage

Для количественной характеристики эффективности методов была
построена диаграмма числа выигранных точек (рис.~8).
Выигранной считается точка $(x_0,x_1)$, в которой соответствующий метод достигает
решения за наименьшее число итераций.

\vspace{10pt}
\setcounter{figure}{7}
\begin{figure}[H]
\centering
\includegraphics[width=0.8\textwidth]{bar_chart.png}
\caption{Число выигранных точек}
\end{figure}


Таким образом, результаты, полученные на сетке начальных данных,
в целом согласуются с выводами, сделанными на основе частного примера,
приведённого в таблице~\ref{tab:tau_compare}.
Наибольшее число выигранных точек демонстрирует метод Modified Secant (MSe),
который существенно превосходит остальные методы.
Высокую эффективность также показывают методы Halley и Ridders,
занимающие второе и третье места.
Метод FP--MSe, являющийся одним из лучших в частном примере,
в общем случае уступает указанным методам, хотя сохраняет достаточно высокую эффективность.
Следовательно, наиболее устойчивыми к выбору начальных приближений
оказались методы MSe, Halley и Ridders.


\newpage

\begin{thebibliography}{9}

\bibitem{Eustafieva}
\href{https://www.mathnet.ru/php/archive.phtml?wshow=paper&jrnid=mzm&paperid=14020&option_lang=rus}{
Евстафьева В.~В.
Об одном типе колебательных решений неавтономной системы
с релейным гистерезисом.
Математические заметки, 2024.}

\bibitem{Kalitkin}{
    Калиткин Н. Н. Численные методы: учеб. пособие. — 2-е изд., исправленное. — СПб.: БХВ-Петербург, 2011. — 592 с.
}


\bibitem{Press}
\href{https://faculty.kfupm.edu.sa/phys/aanaqvi/Numerical%20Recipes-The%20Art%20of%20Scientific%20Computing%203rd%20Edition%20(Press%20et%20al).pdf}{
Press W.~H., Teukolsky S.~A., Vetterling W.~T., Flannery B.~P.
Numerical Recipes: The Art of Scientific Computing.
Cambridge University Press, 2007.}

\bibitem{Badr}
\href{https://www.researchgate.net/publication/360883886_Novel_hybrid_algorithms_for_root_determining_using_advantages_of_open_methods_and_bracketing_methods}{
Badr E., Attiya H., El-Ghamry A.
Novel hybrid algorithms for root determining using advantages of open methods and bracketing methods.
Alexandria Engineering Journal, 2022.}

\bibitem{Steffensen}{
    Steffensen, I.F. (1933) Remarks on Iteration. Scandinavian Actuarial Journal, 1933, 64-72
}

\bibitem{Traub}{
    Traub, J.F. (1964) Iterative Methods for the Solution of Equations. Prentice-Hall, Englewood Cliffs.
}

\bibitem{Dzunic}
\href{https://www.sciencedirect.com/science/article/pii/S0893965912001395}{
Džunić J., Petković M.
A cubically convergent Steffensen-like method for solving nonlinear equations.
Applied Mathematics Letters, 2012.}

\bibitem{Matlab}
MathWorks.
MATLAB Documentation.
Available at:
\url{https://www.mathworks.com/help/matlab/ref/fzero.html}
\url{https://www.mathworks.com/help/matlab/ref/roots.html}
\url{https://www.mathworks.com/help/optim/ug/fsolve.html}
\url{https://www.mathworks.com/help/symbolic/sym.solve.html}
\url{https://www.mathworks.com/help/symbolic/sym.vpasolve.html}
    

\bibitem{Julia}
Julia Programming Language.
Documentation.
Available at:
\url{https://docs.julialang.org/}
\url{https://github.com/JuliaMath/Roots.jl/}

\bibitem{matplotlib}{
    Matplotlib documentation. Available at: \url{https://matplotlib.org/stable/}
}

\end{thebibliography}
\end{document}
